\chapter{Structured arguments, preferences, and defeat}
\label{sec:p18-structured-argumentation}

An abstract argumentation framework records which arguments defeat which
others. A structured framework also specifies how the arguments are built
and where a challenge reaches them. This chapter develops the two-premise-sort
ASPIC+ construction of Modgil and Prakken's tutorial, including its general
contrary function \cite[Sections 3 and 5]{ModgilPrakken2014ASPICTutorial}.
The captured author manuscript is dated December 16, 2013; the article was
published in 2014. The attack relation, preference convention, and admissible
argument trees below are fixed together. Every finite argument tree is
available; no additional consistent-premise or minimal-support restriction
is imposed.

\section{The language and the argument trees}

\begin{hbdefinition}{def:p18-aspic-system}{Typed argumentation data}
An argumentation system consists of a set of formulas \(L\), a function
\(C:L\longrightarrow\mathcal P(L)\), disjoint sets \(R_s,R_d\) of strict and
defeasible rules, and a partial naming function \(n:R_d\rightharpoonup L\).
A rule \(r\) has a finite ordered list
\(\operatorname{ant}(r)=(\phi_1,\ldots,\phi_k)\in L^k\), where \(k\geq0\),
and a consequent \(\operatorname{head}(r)\in L\). We write its data as
\((\phi_1,\ldots,\phi_k\to\psi)\in R_s\) or
\((\phi_1,\ldots,\phi_k\Rightarrow\psi)\in R_d\).
These arrows classify rules; they are not formulas of \(L\).
A rule is identified by its type, antecedent list, and consequent;
repeating the same rule description does not create a second rule.

The knowledge base is a disjoint union \(K=K_n\mathbin{\dot\cup}K_p\subseteq L\).
Members of \(K_n\) are axiom premises; members of \(K_p\) are ordinary,
challengeable premises. Rule names are formulas, not additional premises.
An unnamed defeasible rule has no name-based undercutter. Injectivity of
\(n\) is not required: a challenge to a shared name can reach each rule
with that name.

The statement \(\alpha\in C(\beta)\) says that \(\alpha\) conflicts with the
target \(\beta\). It is a \emph{contrary} conflict if
\(\beta\notin C(\alpha)\), and a \emph{contradictory} conflict if
\(\beta\in C(\alpha)\). Every formula is required to have at least one
contradictory. A contradictory need not be unique. The construction does
not assume that \(C\) is symmetric or that self-conflict is impossible.
The usual paired literals \(p,\neg p\), with \(C(p)=\{\neg p\}\) and
\(C(\neg p)=\{p\}\), give a special case.
\end{hbdefinition}

\begin{hbdefinition}{def:p18-aspic-arguments}{Finite trees and their components}
The set \(\mathcal A\) of arguments is the least set generated as follows.
For each \(\phi\in K\), there is a premise argument \([\phi]\).
If \(r\in R_s\cup R_d\) has antecedents
\((\phi_1,\ldots,\phi_k)\), and arguments \(A_i\) conclude \(\phi_i\),
there is an argument \([A_1,\ldots,A_k;r]\) concluding
\(\operatorname{head}(r)\). A nullary rule gives the distinct argument
\([\,;r]\). Arguments are finite labelled ordered trees, rather than
their conclusions or their sets of premises.

For a premise argument \(A=[\phi]\), put
\[
 \operatorname{Conc}(A)=\phi,\quad
 \operatorname{Prem}(A)=\{\phi\},\quad
 \operatorname{Sub}(A)=\{A\},\quad
 \operatorname{Def}(A)=\operatorname{Last}(A)=\varnothing .
\]
Its top rule \(\operatorname{Top}(A)\) is undefined.
For \(A=[A_1,\ldots,A_k;r]\), put
\[
\begin{aligned}
 \operatorname{Conc}(A)&=\operatorname{head}(r),&
 \operatorname{Top}(A)&=r,\\
 \operatorname{Prem}(A)&=\bigcup_i\operatorname{Prem}(A_i),&
 \operatorname{Sub}(A)&=\{A\}\cup\bigcup_i\operatorname{Sub}(A_i),\\
 \operatorname{Def}(A)&=\bigcup_i\operatorname{Def}(A_i)
                   \ \cup\ \bigl(\{r\}\cap R_d\bigr).
\end{aligned}
\]
If \(r\in R_d\), set \(\operatorname{Last}(A)=\{r\}\); if \(r\in R_s\),
set \(\operatorname{Last}(A)=\bigcup_i\operatorname{Last}(A_i)\).
All empty unions are empty.

Write \(\operatorname{Prem}_p(A)=\operatorname{Prem}(A)\cap K_p\).
An argument is \emph{strict} if \(\operatorname{Def}(A)=\varnothing\),
and otherwise \emph{defeasible}; it is \emph{firm} if
\(\operatorname{Prem}_p(A)=\varnothing\), and otherwise \emph{plausible}.
A strict top rule does not erase defeasible inferences below it.
\end{hbdefinition}

\begin{hblemma}{thm:p18-aspic-tree-properties}{proved=annex:p18.aspictree}{Tree invariants}
The construction gives a set of finite arguments. Every argument has finitely
many subarguments, premises, and defeasible rules.
If \(B\in\operatorname{Sub}(A)\), then
\(\operatorname{Sub}(B)\subseteq\operatorname{Sub}(A)\).
Moreover,
\(\operatorname{Last}(A)=\varnothing\) if and only if \(A\) is strict.
\end{hblemma}

\section{An explicit preference convention}

The relation \(A\prec B\) means that \(A\) is strictly less preferred
than \(B\). It can be supplied directly as an asymmetric relation, including
the empty relation. The following alternatives compute it from two
independent preorders: \(\leq_d\) on \(R_d\), and \(\leq_p\) on \(K_p\).
Neither preorder must be total. An axiom is outside the domain of
\(\leq_p\), so every premise comparison below uses
\(\operatorname{Prem}_p\).
For \(i\in\{p,d\}\), write \(x<_i y\) exactly when
\(x\leq_i y\) and \(\neg(y\leq_i x)\).

\begin{hbdefinition}{def:p18-aspic-set-lifting}{Finite-set comparisons}
For a preorder \(\leq\) on a set \(X\), define relations
\(\unlhd_{\rm E},\unlhd_{\rm D}\) on finite subsets of \(X\).
For either relation, an empty left operand never compares below another
operand; a nonempty left operand compares below an empty right operand.
For nonempty \(\Gamma,\Delta\), define
\[
\begin{split}
 \Gamma\unlhd_{\rm E}\Delta
   &\iff (\exists x\in\Gamma)(\forall y\in\Delta)\ x\leq y,\\
 \Gamma\unlhd_{\rm D}\Delta
   &\iff (\forall x\in\Gamma)(\exists y\in\Delta)\ x\leq y .
\end{split}
\]
The subscripts denote the elitist and democratic comparisons, respectively.
For either chosen comparison, its strict part is
\[
 \Gamma\lhd\Delta
 \iff \Gamma\unlhd\Delta\ \hbox{and}\ \neg(\Delta\unlhd\Gamma).
\]
The empty-set clauses apply before the quantified clauses.
The strict part is calculated from the complete relation; it is not defined
by replacing each occurrence of \(\leq\) with its strict part.
\end{hbdefinition}

\begin{hbtheorem}{thm:p18-aspic-set-properties}{proved=annex:p18.aspicsets}{Properties and a strict-comparison boundary}
Both set comparisons are transitive, and their strict parts are irreflexive
and transitive. The set comparisons need not be reflexive:
in particular, \(\varnothing\not\unlhd\varnothing\).
On nonempty finite subsets of a totally preordered set, the elitist
comparison is comparison of a least element of each set; the democratic
comparison is comparison of a greatest element.

For partial preorders, replacing every base weak comparison by a strict
comparison can change the strict set relation. With incomparable
\(a,b\), the sets \(\{a,b\}\) and \(\{a\}\) witness this for the elitist
comparison.
\end{hbtheorem}

\begin{hbdefinition}{def:p18-aspic-argument-lifting}{Last-link and weakest-link comparisons}
Choose one of the two set comparisons, using \(\leq_p\) for premise sets
and \(\leq_d\) for rule sets. These typed instances are both denoted
\(\unlhd\).
Define the last-link comparison by
\[
\begin{split}
A\preceq_{\rm L}B\iff{}&
 \operatorname{Last}(A)\unlhd\operatorname{Last}(B)\\
&{}\lor\bigl(
 \operatorname{Last}(A)=\operatorname{Last}(B)=\varnothing
 \ \land\
 \operatorname{Prem}_p(A)\unlhd\operatorname{Prem}_p(B)\bigr).
\end{split}
\]
Define the weakest-link comparison by the following ordered cases:
\[
 A\preceq_{\rm W}B\iff
 \begin{cases}
 \operatorname{Prem}_p(A)\unlhd\operatorname{Prem}_p(B),
       &A,B\text{ both strict},\\
 \operatorname{Def}(A)\unlhd\operatorname{Def}(B),
       &A,B\text{ both firm and not both strict},\\
 \operatorname{Prem}_p(A)\unlhd\operatorname{Prem}_p(B)
   \ \land\ \operatorname{Def}(A)\unlhd\operatorname{Def}(B),
       &\text{otherwise}.
 \end{cases}
\]
In either case, use exactly
\[
 A\prec_{\rm L}B\iff
 A\preceq_{\rm L}B\land\neg(B\preceq_{\rm L}A),
 \qquad
 A\prec_{\rm W}B\iff
 A\preceq_{\rm W}B\land\neg(B\preceq_{\rm W}A).
\]
Thus the resulting strict relation is asymmetric. This definition makes no
additional claim that every such argument comparison has all properties
required by a general consistency theorem.
\end{hbdefinition}

\section{Where attacks land and when they defeat}

\begin{hbdefinition}{def:p18-aspic-attacks}{Three attacks at identified subarguments}
Let \(A,B\in\mathcal A\). An attack includes its target occurrence type.
The following conditions define the three possibilities.

\emph{Undermining:} a premise subargument \(B'=[\phi]\) of \(B\) has
\(\phi\in K_p\), and \(\operatorname{Conc}(A)\in C(\phi)\).

\emph{Rebutting:} a subargument \(B'\) of \(B\) has a top rule in \(R_d\),
and \(\operatorname{Conc}(A)\in C(\operatorname{Conc}(B'))\).

\emph{Undercutting:} a subargument \(B'\) of \(B\) has top rule
\(r\in R_d\cap\operatorname{dom}(n)\), and
\(\operatorname{Conc}(A)\in C(n(r))\).

An undermining or rebutting attack is called contrary when its displayed
conflict is asymmetric, and contradictory when that conflict is symmetric.
An argument may attack another in more than one way. Axiom premises and
strict top-rule applications are not undermining or rebutting targets.
Their superarguments may nevertheless contain other, attackable targets.
\end{hbdefinition}

\begin{hbdefinition}{def:p18-aspic-defeat}{Preference at the point of attack}
Fix one asymmetric strict preference relation \(\prec\) on \(\mathcal A\).
An undercut always succeeds. A contrary undermining or contrary rebutting
attack always succeeds. A contradictory undermining attack at
\(B'=[\phi]\) succeeds exactly when \(\neg(A\prec B')\).
A contradictory rebutting attack at \(B'\) succeeds exactly when
\(\neg(A\prec B')\).

Write \(A\,D\,B\) if at least one attack of \(A\) on \(B\) succeeds.
The preference comparison is against the attacked subargument \(B'\);
it is not a comparison against the whole superargument \(B\).
Equality of strength and incomparability both satisfy
\(\neg(A\prec B')\).
\end{hbdefinition}

\begin{hblemma}{thm:p18-aspic-locality}{proved=annex:p18.aspiclocality}{Propagation and strict roots}
If \(B\in\operatorname{Sub}(T)\), every attack and every defeat of \(B\)
also attacks or defeats \(T\), respectively.
If \(T=[A_1,\ldots,A_k;r]\) with \(r\in R_s\), every defeater of \(T\)
defeats at least one \(A_i\). Consequently a nullary strict argument has
no defeaters.
\end{hblemma}

\begin{hbdefinition}{def:p18-aspic-semantics}{Translation and conclusion queries}
The induced abstract framework is \(F_{\rm AS}=(\mathcal A,D)\).
Use the complete, grounded, preferred, and stable semantics of
Definitions~\ref{def:p18-extension-semantics} and
\ref{def:p18-arbitrary-framework}, with \(D\) as the edge relation for
both conflict-freeness and defense. Structural attacks that fail the
preference condition are not edges of this framework.

For an extension \(E\), set
\(\operatorname{Conc}(E)=\{\operatorname{Conc}(A):A\in E\}\).
For a specified nonempty family \(\mathcal E\) of extensions, distinguish
three formula questions:
\[
\begin{array}{ll}
\text{credulous:}&
  (\exists E\in\mathcal E)(\exists A\in E)\
                  \operatorname{Conc}(A)=\phi,\\
\text{common-argument skeptical:}&
  (\exists A\in\bigcap\mathcal E)\
                  \operatorname{Conc}(A)=\phi,\\
\text{extensionwise skeptical:}&
  (\forall E\in\mathcal E)(\exists A\in E)\
                  \operatorname{Conc}(A)=\phi .
\end{array}
\]
All three are declared false when \(\mathcal E=\varnothing\); the absence
of extensions is also reported explicitly, as in
Definition~\ref{def:p18-acceptance-queries}.
The tutorial's formula-skepticism uses the common-argument reading.
Its exclusive use of \emph{credulous} can be recovered by adding
\emph{and not common-argument skeptical} to the first query.
\end{hbdefinition}

\begin{hbtheorem}{thm:p18-aspic-query-distinction}{proved=annex:p18.aspicqueries}{Different witnesses in different extensions}
Common-argument skepticism implies extensionwise skepticism, but the
converse fails even in a finite ASPIC+ system with only ordinary premises
and defeasible rules. An explicit four-argument example is supplied in
the proof.
\end{hbtheorem}

\section{Closure and the limits of consistency}

\begin{hbdefinition}{def:p18-aspic-consistency}{Strict closure and consistency}
For \(S\subseteq L\), put \(S_0=S\) and
\[
 S_{j+1}=S_j\cup
 \{\operatorname{head}(r):r\in R_s,\
       \text{every antecedent of }r\text{ belongs to }S_j\}.
\]
The strict closure is \(\operatorname{Cl}_s(S)=\bigcup_{j<\omega}S_j\).
A set \(S\) is directly consistent if no
\(\alpha,\beta\in S\) satisfy \(\alpha\in C(\beta)\);
it is indirectly consistent if \(\operatorname{Cl}_s(S)\) is directly
consistent. These definitions allow \(\alpha=\beta\).
For a language of paired literals, a transposition of
\((\phi_1,\ldots,\phi_k\to\psi)\), at position \(i\), is the strict
rule obtained by replacing antecedent \(\phi_i\) with the paired
contradictory \(\neg\psi\) and replacing the consequent with
\(\neg\phi_i\). Here pairing is an involution, so
\(\neg(\neg\phi)=\phi\). A strict-rule set is closed under
transposition if it contains these rules for every rule and every
antecedent position.
\end{hbdefinition}

\begin{hbtheorem}{thm:p18-aspic-structural-rationality}{proved=annex:p18.aspicclosure}{Subarguments and strict consequences in complete extensions}
For every complete extension \(E\) of \(F_{\rm AS}\),
\[
 A\in E\Longrightarrow\operatorname{Sub}(A)\subseteq E,
 \qquad
 \operatorname{Cl}_s(\operatorname{Conc}(E))
       =\operatorname{Conc}(E).
\]
The statement holds with infinite languages and argument sets, arbitrary
finite-arity strict rules, and either preference construction above.
It also holds for grounded, preferred, and stable extensions.
It does not assert that these conclusion sets are consistent.
\end{hbtheorem}

\begin{hbtheorem}{thm:p18-aspic-flat-consistency}{proved=annex:p18.aspicflat}{Consistency when every root is challengeable}
Suppose \(K_n=R_s=\varnothing\) and \(\prec\) is asymmetric.
Every conflict-free set of arguments has a directly consistent set of
conclusions. Consequently every complete extension satisfies both forms
of consistency and both closure properties of the preceding theorem.
\end{hbtheorem}

\begin{hbtheorem}{thm:p18-aspic-rationality-boundaries}{proved=annex:p18.aspicboundaries}{Three failures with explicit finite systems}
Each of the following can occur.
First, conflicting axiom premises belong to the unique extension.
Second, directly consistent ordinary premises and unopposed strict
inferences produce inconsistent conclusions when the strict-rule set
lacks the needed transpositions.
Third, changing rebuttal to permit attacks at strict roots can produce a
stable extension that is not closed under strict inference, even when
the strict rules occur in transposed pairs.
The proof gives all arguments and defeats in each system.
\end{hbtheorem}

\section{A complete competing-argument calculation}

Use paired positive and negative literals for
\(p,q,r,s,t,u,v,x,d_1,\ldots,d_6\), with the symmetric conflict convention
of Definition~\ref{def:p18-aspic-system}. Set \(K_n=\{p\}\),
\(K_p=\{s,u,x\}\). The two strict rules have data
\[
 s_1:(p,q)\to r,\qquad s_2:v\to\neg s,
\]
and the six defeasible rules have data
\[
\begin{array}{lll}
 d_1:p\Rightarrow q,&d_2:s\Rightarrow t,&d_3:t\Rightarrow\neg d_1,\\
 d_4:u\Rightarrow v,&d_5:(v,x)\Rightarrow\neg t,&
 d_6:s\Rightarrow\neg p .
\end{array}
\]
Set \(n(d_i)=d_i\); the symbols on the right denote the corresponding
name formulas. All arguments are listed below. This reconstruction
includes the argument generated by \(d_6\).
\[
\begin{array}{c|l|c|c}
\text{argument}&\text{tree}&\operatorname{Prem}_p&\operatorname{Def}\\ \hline
A_1&[p]&\varnothing&\varnothing\\
A_2&[A_1;d_1]&\varnothing&\{d_1\}\\
A_3&[A_1,A_2;s_1]&\varnothing&\{d_1\}\\
B_1&[s]&\{s\}&\varnothing\\
B_2&[B_1;d_2]&\{s\}&\{d_2\}\\
B_3&[B_2;d_3]&\{s\}&\{d_2,d_3\}\\
B_4&[B_1;d_6]&\{s\}&\{d_6\}\\
C_1&[u]&\{u\}&\varnothing\\
C_2&[C_1;d_4]&\{u\}&\{d_4\}\\
C_3&[C_2;s_2]&\{u\}&\{d_4\}\\
D_3&[x]&\{x\}&\varnothing\\
D_4&[C_2,D_3;d_5]&\{u,x\}&\{d_4,d_5\}
\end{array}
\]
In particular, \(A_3,C_3\) have strict top rules but are defeasible
arguments. The top rule of \(D_4\) is defeasible.

\begin{hbtheorem}{thm:p18-aspic-worked-graph}{proved=annex:p18.aspicgraph}{All attacks and the unranked extensions}
The displayed list is exhaustive. All structural attack edges, grouped
by their source and target sets, are
\[
\begin{array}{c|c|l}
\text{source}&\text{targets}&\text{mode}\\ \hline
B_3&\{A_2,A_3\}&\text{undercut at }A_2\\
C_3&\{B_1,B_2,B_3,B_4\}&\text{undermine at }B_1\\
D_4&\{B_2,B_3\}&\text{rebut at }B_2\\
B_2&\{D_4\}&\text{rebut at }D_4\\
A_1&\{B_4\}&\text{rebut at }B_4 .
\end{array}
\]
There are no further edges. In particular \(B_4\), despite concluding
\(\neg p\), does not attack the axiom argument \(A_1\).

When \(\prec=\varnothing\), every displayed attack is a defeat.
There is a unique complete, grounded, preferred, and stable extension,
\[
 E_0=\{A_1,A_2,A_3,C_1,C_2,C_3,D_3,D_4\}.
\]
The conclusion \(r\) is therefore skeptically accepted under either formula
reading for each of these semantics.
\end{hbtheorem}

\begin{hbtheorem}{thm:p18-aspic-worked-preferences}{proved=annex:p18.aspicpreferences}{The two link principles give different conclusions}
For the same system, use the elitist set comparison and priorities
\(u<_p s,\ x<_p s,\ d_4<_d d_2<_d d_5\).
One concrete total-preorder choice assigns ranks \(0,0,1\) to \(u,x,s\),
respectively, and ranks \(0,1,2\) to \(d_4,d_2,d_5\);
assign any ranks to \(d_1,d_3,d_6\). Lower rank means less preferred.

Under last-link preferences the unique extension in all four semantics is
\[
 E_{\rm L}=\mathcal A\setminus\{B_2,B_3,B_4\}.
\]
Under weakest-link preferences it is
\[
 E_{\rm W}=\mathcal A\setminus\{A_2,A_3,B_4,D_4\}.
\]
Thus \(r\) is accepted in the first case and is absent from every extension
in the second. Both extensions contain \(s,\neg s\)
as conclusions; these data therefore do not satisfy a general consistency
guarantee.
\end{hbtheorem}

\begin{hbtheorem}{thm:p18-aspic-contrary-example}{proved=annex:p18.aspiccontrary}{An asymmetric contrary defeats a stronger premise}
Let \(L=\{p,\bar p,s,\bar s\}\), with
\[
 C(p)=\{\bar p\},\quad C(\bar p)=\{p\},\quad
 C(s)=\{\bar s,p\},\quad C(\bar s)=\{s\}.
\]
Take \(K_p=\{p,s\}\), \(K_n=R_s=R_d=\varnothing\), and \(p<_p s\).
For either link comparison, \([p]\prec[s]\), but \([p]\) defeats \([s]\).
The unique extension in all four semantics is \(\{[p]\}\).
Every formula in this example has a contradictory, although the additional
conflict from \(p\) to \(s\) is asymmetric.
\end{hbtheorem}

\section{When the construction is a finite decision procedure}

\begin{hbtheorem}{thm:p18-aspic-finite-generation}{proved=annex:p18.aspicfinite}{Finite generation under an acyclic dependency rank}
Suppose \(K,R_s,R_d\) are finite and every formula occurring in these data
has a rank in \(\{0,\ldots,H\}\), for a nonnegative integer \(H\),
such that every antecedent of every
non-nullary rule has smaller rank than its consequent.
Then \(\mathcal A\) is finite. Generate premise and nullary-rule arguments,
and then process rule consequents in increasing rank, forming every ordered
tuple of already generated arguments with the required conclusions.
This procedure generates each argument tree and terminates.

If conflict membership, rule names, and the chosen priorities on these finite
data are explicitly given, the attack and defeat relations are decidable by
the displayed definitions. The existing finite extension methods then decide
argument and formula acceptance. Finite \(K,R_s,R_d\) alone are insufficient:
with \(K_n=\{p\}\) and \(R_s=\{p\to p\}\), there are infinitely many arguments.
\end{hbtheorem}

\begin{proofboundary}{def:p18-aspic-next-boundary}{Further source results}
The construction and the proved closure results above are available without
a general consistency theorem. The next stronger source result concerns
axiom consistency, a precise contraposition or transposition condition,
and the technical reasonable-order condition of the 2013 general account.
For asymmetric contraries its additional well-formedness condition excludes
a target of a one-way contrary from both \(K_n\) and the consequents of
\(R_s\). These conditions must be stated and their complete derivation
supplied before invoking that stronger theorem.

Separate further constructions include the tutorial's classical and Tarskian
instantiations with restricted argument supports, its assumption-based
translation, and its detailed argument-scheme models. Their additional
restrictions are not implicit in the all-finite-trees construction here.
\end{proofboundary}
